\section{I/O performance}
\label{section:tailoring:io}

% again, there is not a huge amount we can test here. If the code stipulates any different I/O patterns then we can try these. Otherwise, on a particular system, we can try changing where files are read from and written to, e.g., a temporary storage on the node.

If a code performs heavy I/O, and the high-level metrics demonstrate this to be
a potential bottleneck, then we must check whether we are capturing realistic
I/O performance before proceeding. It is easy to poorly use the available file
systems on a machine resulting in excessive wait times. To be more explicit:
are we using our system well enough to see the I/O perform ``in its best
light"? 

Many codes come with build and run instructions that are agnostic about where
the source code, binaries and any data files are located. This makes sense as
the instructions are intended to be flexible to be used on a variety of
systems. However, for codes that are heavily I/O bound, we are not best advised
to read and write files from home directories or equivalent. Instead, we will
find better performance by doing our I/O on the node-local file system.
Typically, file systems like the Network File System (NFS) will have high
latency and low bandwidth, whereas node-local file systems --- usually SSDs on
the node --- have low latency and high bandwidth which we can exploit for our
I/O heavy workloads.

Often node-local file systems are labelled as temperary directories (e.g.,
\texttt{\$TEMPDIR}, \texttt{\$TMPDIR}, \texttt{\$TMP}, \ldots), and so a simple
workflow of this so-called \emph{data staging} can be defined as:
\begin{enumerate}
  \item Copy data files or unpack \texttt{tar}, etc. files to the temporary directory
  \item Run code, ensuring to set data paths to the temporary directory
  \item Move any necessary output files from the temporary directory back to the NFS as they will otherwise be lost after job completion
\end{enumerate}
These workflows are particularly relevant to codes that rely on I/O of many
small files, e.g., ML codes.

The difficulty with tailoring the I/O of a code on a particular system, is that
these features can vary system to system due to hardware or the way that the
system is configured. Many compute centres give their users information on file
systems and relative performance, so we always recommend consulting
documentation for any other available features of the file system we can
exploit.
